
%% Abstract text
%% All the details (name, title, etc.) on the abstract page appear as specified
%% above. Add your abstarct text with paragraphs here to have paragraphs in the
%% visible abstract page. Nonetheless, write the abstarct text without
%% paragraphs in the macro \thesismacro so that it is added to the metadata as
%% well.
%%
%% Abstract keywords
\keywords{Retrieval-Augmented Generation, RAG, prompt caching, model cascading, semantic re-ranking, enterprise systems}
\begin{abstractpage}[english]
Retrieval augmented generation is widely used method for enriching generative AI applications in enterprise context. It has couple of major disadvantages, such as hallucination and high costs. This thesis examines four distinct methods to optimize the RAG pipeline and those methods are measured against optimal combination of latency, costs and quality. Studied methods are Prompt caching, semantic re-ranking, small-to-big retireval and model cascading. 

The main contribution is to test all of the tested both, in isolation, as well as in combination using enterprise ready frameworks such as Java and Spring AI. In addition, tests were ran using commercial Large Language Model API endpoints in cloud platform. All the optimization methods were ran against two distinct dataset. The first dataset is an academic benchmark for RAG applications, and the second one is curated enterprise customer support system dataset.

The main cost reduction driver is prompt caching, which provides exceptional cost benefits both in isolation as well as a part of the combination. Interestingly, none of the studied methods reduced latency, but had only increasing impact. Pareto optimal combination on both quality per dollar and quality per second for both datasets is Small-to-Big retrieval combined with the prompt caching. Small-to-big drives the quality up, while the prompt caching reduces the costs.

If only the costs are the only prompt caching combined with the model cascading is the most optimal strategy for enterprise dataset. Mainly, model cascading works best on enterprise dataset, where its cheaper variants can be used more often, while on more complex academic benchmark, it acted as a cost increase. The semantic re-ranking did not provide any meaningful benefits on quality per dollar or second. The findings provides useful indications for enterprise architects, when optimizing generative AI applications in enterprise systems.  
\end{abstractpage}

%% The text in the \thesisabstract macro is stored in the macro \abstractext, so
%% you can use the text metadata abstract directly as follows:
%%
%\begin{abstractpage}[english]
%	\abstracttext{}
%\end{abstractpage}

%% Force a new page so that the possible Finnish or Swedish abstract does not
%% begin on the same page
%%
\newpage
%%
%% Abstract in Finnish. Delete if you don't need it. 
%%
%% Respecify those fields that differ from the earlier specification or simply
%% respecify all fields.
\thesistitle{Context is king - Retrieval-Augmented Generation in enterprise-ready systems}
\supervisor{Prof. Petri Vuorimaa}
\advisor{DI Lauri Salminen}
\degreeprogram{Computer, Communication and Information Sciences}
\major{Computer Science}
\collaborativepartner{Nordea Bank Oyj}
\date{30.6.2026}
%% The keywords need not be separated by \spc now.
\keywords{hakuavusteinen generointi, RAG, kehotteen välimuisti, mallireititys, semanttinen uudelleenjärjestäminen, yritysjärjestelmät}
%% Abstract text
\begin{abstractpage}[finnish]
Tiivistelmä on lyhyt kuvaus työn keskeisestä sisällöstä usein yhtenä kappaleena:
mitä tutkittiin ja miten sekä mitkä olivat tärkeimmät tulokset. Suomenkielisen
opinnäytteen tiivistelmä kirjoitetaan suomeksi ja englanniksi ja ruotsinkielisen
vastaavasti ruotsiksi ja englanniksi. Suomen- tai ruotsinkielisten
opiskelijoiden, joiden opinnäytteen kieli on englanti, tulee kirjoittaa
tiivistelmänsä englanniksi ja koulusivistyskielellään. Muiden kuin
koulusivistyskieleltään suomen- tai ruotsinkielisten tulee kirjoittaa
tiivistelmänsä vain englanniksi. Opiskelija voi halutessaan lisätä
opinnäytteeseensä toisen tai kolmannen tiivistelmän omalla äidinkielellään.
Tämän opinnäytteen tiivistelmäteksti kirjoitetaan opinnäytteen luettavan osan
lomakkeen lisäksi myös pdf-tiedoston metadataan. Kirjoita tähän metadataan
kirjoitettavaa teksti. Metadatatekstissa ei saa olla erikoismerkkejä,
rivinvaiho- tai kappaleenjakomerkkiä, joten näitä merkkeja ei saa käyttää tässä.
Jos tiivistelmäsi ei sisällä erikoimerkkejä eikä kaipaa kappaleenjakoa, voit
hyödynttää makroa abstracttext luodessasi lomakkeen tiivistelmää (katso
kommentti tässä TeX-tiedostossa alla). Metadatatiivistelmatekstin on muuten 
oltava sama kuin lomakkeessa oleva teksti.

\end{abstractpage}
